% P51 | 6 Comparison of methods | Result: the European call on one footing
% guide: dev_notes/paragraph_guides/P51_result_the_european_call_on_one_footing.txt
Table~\ref{tbl:deskcall} shows that the vectorised CRR tree with $N=252$ reproduces the price of the list-based class, with an error of \DeskCrrVecErr, in \DeskCrrVecMs\ ms instead of \DeskCrrClassMs\ ms.
The Leisen--Reimer tree with $N=101$ has an error of \$\DeskLrErr\ in \DeskLrMs\ ms, and CRR with smoothing and Richardson extrapolation at $N=252$ has an error of \$\DeskSmoothErr.
Plain Monte Carlo with $10^5$ samples has a standard error of \$\DeskMcPlainSE\ (the realised error is \$\DeskMcPlainErr) and takes \DeskMcPlainMs\ ms.
The control variate brings the standard error to \$\DeskMcCvSE\ at the same cost, and scrambled Sobol points with $2^{16}$ samples reach an error of \$\DeskSobolErr\ in \DeskSobolMs\ ms.
The Euler--Maruyama simulation with $252$ steps takes about \DeskMcEulerMs\ ms for the same $10^5$ samples, roughly two orders of magnitude more than the exact sampler, and has the same standard error.
For a European payoff on geometric Brownian motion the extra cost buys nothing, because the exact terminal law is known; the scheme earns its place when the dynamics are not solvable.
The order of the rows is stable across machines, but the absolute times are not and should be read as orders of magnitude.
